\subsection{The Feedback Loop: Synchronic Chain, Diachronic Saturation}

The closing thesis of \S1.6 --- that emotion is the soul's kinetic alignment toward a world disclosed as mattering in a particular way --- leaves a structural feature implicit in the \textit{Rhetoric}'s judgment-affecting definition unaddressed. The linear chain $A_0 \rightarrow A_4$ established in the prior section captures the synchronic transitions through which basic valence grounds \textit{doxa}, \textit{doxa} concretizes the higher-order \textit{path\=e}, and the higher-order \textit{path\=e} mobilize action; yet emotion, once produced, alters the very conditions under which subsequent \textit{doxai} will be formed. Specifically, the chain is not a closed circuit describable through forward transitions alone: each actualization persists as dispositional ground for what follows. This section names that persistence the \textit{pathos} feedback loop and analyzes its structure on two registers, the synchronic and the diachronic, before grounding the diachronic register in the temporality of \textit{pathos}-grade thrownness.

The \textit{locus classicus} for the feedback's articulation in Aristotle's psychological corpus is \textit{De Insomniis}. Aristotle there observes that ``we are easily deceived respecting the operations of sense-perception when we are excited by emotions, and different persons according to their different emotions; for example, the coward when excited by fear, the amorous person by the object of his desire; so that, with but little resemblance to go upon, the former thinks he sees his foes approaching, the latter, that he sees the object of his desire; and the more deeply one is under the influence of the emotion, the less similarity is required to give rise to these impressions'' (Aristotle, \textit{De Insomniis} 460b3--11). The passage isolates a specific cognitive mechanism: under the sway of an active emotion, the controlling sense --- the corrective faculty that in ordinary waking cognition intercepts misleading \textit{phantasmata} before they can issue in \textit{doxa} --- is impaired, and \textit{phantasmata} that would otherwise be filtered are instead taken as true. Aristotle traces the impairment, indeed, to a structural fact of the cognitive apparatus: the faculty by which the controlling sense judges is not identical with the faculty by which images come before the mind, and the latter can override the former when the perceiving organism is sufficiently agitated.

The result is a self-reinforcing recursion. An initial \textit{doxa} produces an initial \textit{pathos}; the \textit{pathos} impairs the corrective faculty; subsequent \textit{phantasmata} pass through the corrective gate into \textit{doxa} unopposed; the new \textit{doxai} produce further \textit{path\=e}; and the cycle continues until external intervention, somatic exhaustion, or recovery of the corrective faculty breaks the recursion. The mechanism Aristotle describes is, accordingly, sufficiently general to render intelligible a range of phenomena that are otherwise puzzling. The escalation of anger through sustained rumination on a perceived slight, the intensification of fear through repeated imagining of an anticipated evil, the susceptibility of the dreamer, the ill, and the passionately agitated to phantasmatic deception (Aristotle, \textit{De Insomniis} 460b3--16; 461b3--8) --- each reflects the same loop running unchecked in conditions where the corrective faculty is structurally weakened. Similarly, the rhetor's strategy that the \textit{Rhetoric} II.1 makes explicit becomes intelligible from this side: by first establishing an emotional attunement in the audience --- pity for the client, indignation toward the opponent, fear of the projected consequence --- the rhetor weakens the audience's corrective faculties and thereby renders them more susceptible to the \textit{phantasmata} that vivid language (ἐνέργεια, \textit{energeia}) presents to their imaginations.

Two structural features of the mechanism deserve emphasis. First, the loop is not a malfunction of the cognitive system but an intrinsic feature of the relationship between emotion and \textit{doxa} in rational animals. Aristotle's claim is not that the corrective faculty fails because of a defect to be remedied but that the architecture of judgment is such that emotional intensification produces, of itself, a corresponding lowering of the threshold at which \textit{phantasmata} are admitted into \textit{doxa}. Second, the recursion is not symmetric in time: an initial \textit{pathos} lowers the threshold for the \textit{next} \textit{doxa}, which in turn produces a \textit{pathos} whose intensity reflects the lowered threshold. Thus the temporal direction of the modulation runs from past to present: prior \textit{path\=e} modulate present \textit{doxa}-formation, and present \textit{doxa}-formation feeds forward into present \textit{path\=e}. This asymmetry will become decisive when, in the analysis to follow, the feedback is shown to belong to the diachronic rather than the synchronic register of the chain.\footnote{The feedback loop, as developed here, is an interpretive reconstruction rather than a doctrine Aristotle states in these terms. Three passages converge on the structure: \textit{De Anima} III.3, 427b21--24 establishes that \textit{doxa} produces emotion immediately; \textit{De Insomniis} 460b3--16 establishes that emotional states permit \textit{phantasmata} to be taken as true when the controlling faculty is impaired; \textit{Rhetoric} II.1, 1378a20--22 establishes that the \textit{path\=e} modify \textit{kriseis}. That these passages describe aspects of a single recursive mechanism is a plausible but not indisputable interpretive claim; each passage addresses, strictly speaking, a different phenomenon. The synthesis is, however, philosophically productive, and the linear $A_0 \rightarrow A_4$ notation of \S1.6's chain cannot fully represent it: a complete topological representation would require a return path from the \textit{pathos} node at $A_3$ back to the gating condition within the $A_2 \rightarrow A_3$ transition.}

Two registers are at work in the feedback phenomenon, and they must be distinguished if the loop's structural function is not to be confused with the chain's directional architecture. The synchronic register concerns the structural transitions at any single moment of \textit{pathos}-actualization: basic valence at $A_0$ grounds \textit{doxa} at $A_1$; \textit{doxa} concretizes the higher-order \textit{path\=e} at $A_2$ to $A_3$; the higher-order \textit{path\=e} mobilize action at $A_4$. These transitions are directional and irreversible within a given chain-actualization. The diachronic register, by contrast, concerns how prior actualizations of the chain modulate subsequent actualizations: a \textit{pathos}-actualization at $T_1$ persists, in the disposition-ground of the affected organism, as a condition shaping the \textit{pathos}-actualization at $T_2$. Accordingly, the \textit{De Insomniis} phenomenon belongs to the diachronic register, not to the synchronic register. The chain remains structurally directional from $A_0$ to $A_4$ within any single actualization; the ``loop'' is the temporal persistence of prior actualizations as ground for subsequent actualizations, not a third synchronic transition appended to the chain.

Two distinct mechanisms operate within this diachronic register, and conflating them obscures the precise locus at which prior \textit{path\=e} exert their force. The first mechanism --- \textit{Stimmung}-saturation --- modulates the \textit{input} to \textit{doxa}. A sustained higher-order \textit{pathos} persists as a tonal coloring that biases subsequent perception; the angry person, indeed, sees enemies in shadows because the anger-\textit{Stimmung} re-saturates the perceptual encounter, so that the basic valence delivered to \textit{doxa} is already shaped by the residue of prior actualizations. The second mechanism --- corrective-faculty impairment, the mechanism Aristotle specifies in \textit{De Insomniis} --- modulates the \textit{gating} of \textit{doxa}: emotional excitement disables the controlling sense, so that \textit{phantasmata} which would normally be intercepted pass through the corrective gate into \textit{doxa} unopposed. The two mechanisms are commonly conflated under the single rubric of the feedback loop, but they operate at different points of the chain: mechanism (1) at the basic-valence input, mechanism (2) at the corrective gate within the $A_2 \rightarrow A_3$ transition. Both are diachronic modulations of synchronic transitions, not new synchronic transitions; to specify that the basic valence at $A_0$ is itself a function of prior actualizations is not to add a transition $A_{-1} \rightarrow A_0$ to the chain but to specify that the \textit{content} of $A_0$ on any given pass is conditioned by the residues of prior passes.

This temporal-dispositional persistence is what Aristotle elsewhere analyzes under the doctrine of ἕξις (\textit{hexis}, ``settled disposition'') and what Heidegger names \textit{Geworfenheit}, thrownness; both name a structure in which prior dispositional actualizations persist as the ground from which present disclosure takes its character. The Aristotelian articulation begins with the analysis of habituation in \textit{Nicomachean Ethics} II.1, where Aristotle states that ``moral excellence comes about as a result of habit'' (Aristotle, \textit{NE} II.1, 1103a16--17). The claim's significance for the present analysis lies in its specification of \textit{how} the disposition is acquired: through repeated actualization of the chain itself. As Aristotle puts the point, ``for the things we have to learn before we can do, we learn by doing, e.g. men become builders by building and lyre-players by playing the lyre; so too we become just by doing just acts, temperate by doing temperate acts, brave by doing brave acts'' (Aristotle, \textit{NE} II.1, 1103a31--b2). Each actualization of a \textit{pathos}-chain leaves a residual modification of the dispositional ground; over time, repeated actualizations consolidate into \textit{hexeis} that condition subsequent chain-actualizations. \textit{Hexis} is therefore not a third stage of the chain alongside basic valence, \textit{doxa}, and the higher-order \textit{path\=e}; it is the \textit{temporal sedimentation} of repeated chain-actualizations into the dispositional ground from which subsequent actualizations will draw their basic valence and within which subsequent \textit{doxai} will be gated.

The Heideggerian articulation runs structurally parallel. Heidegger specifies in \textit{Being and Time} \S29 that state-of-mind (\textit{Befindlichkeit}) discloses Dasein in three essential respects: it discloses Dasein in its thrownness; it has already disclosed Being-in-the-world as a whole, thereby making it possible to direct oneself toward something at all; and it discloses the worldhood of the world in the mode of circumspective concern (SZ \S29, H.137; Eng. 176). Thrownness names, accordingly, the diachronic persistence of prior dispositional states: ``this characteristic of Dasein's Being --- this `that it is' --- is veiled in its `whence' and `whither', yet disclosed in itself all the more unveiledly; we call it the `thrownness' [\textit{Geworfenheit}] of this entity into its `there'; indeed, it is thrown in such a way that, as Being-in-the-world, it is the `there''' (SZ \S29, H.135; Eng. 174). Dasein, in other words, is always already disposed; the disposition is not a contingent psychological state superadded to a pre-existing self but the existential-temporal ground from which any present self-relation takes its bearings.

Heidegger's decisive temporal claim, however, comes not in \S29 but in \S68b, where he specifies that the \textit{Stimmung} itself temporalizes primarily in the past-modality of \textit{Gewesenheit}: ``Understanding is grounded primarily in the future; one's \textit{state-of-mind}, however, temporalizes itself \textit{primarily} in \textit{having been}. Moods temporalize themselves --- that is, their specific ecstasis belongs to a future and a Present in such a way, indeed, that these equiprimordial ecstases are modified by having-been'' (SZ \S68b, H.340; Eng. 390). Thus the having-already-found-itself-disposed of \textit{Befindlichkeit} is not a past \textit{now stored in memory} but a \textit{Gewesenheit} --- a having-been --- that \textit{currently} conditions how the entities of the world can come close. The thrown disposition is operative \textit{as} the present disclosure's condition, not as a remembered antecedent of it.

The convergence of the Aristotelian and Heideggerian articulations is, accordingly, no mere analogy: it is the recognition of one phenomenon --- prior dispositional actualization persists as the ground from which present disclosure takes its character --- described at different ontological registers. Aristotle's articulation operates at the level of perception-and-affection within the soul-structure; Heidegger's operates at the level of Dasein's fundamental disclosedness. The lecture course of 1924 collected as \textit{Basic Concepts of Aristotelian Philosophy} makes the convergence explicit. There Heidegger derives four basic aspects of being-there from the Aristotelian analysis of \textit{hexis}, and the fourth specifies that ``Being-there must, \textit{for itself}, take the opportunity to cultivate this being-composed as a possibility'' (BCAP, p. 122). \textit{Hexis}, Heidegger insists, can have its genesis only ``in being-there itself'' (BCAP, p. 122). The lectures' formulation, indeed, is the prefiguration of what \textit{Being and Time} will name \textit{Geworfenheit}: the dispositional ground that is neither imposed upon Dasein from without nor chosen by it from within but cultivated, in the Aristotelian register, through the very activity that the disposition itself conditions. The feedback loop, accordingly, is not a structurally distinct third direction added to the chain's two synchronic directions; it is the \textit{thrownness-character} of the chain --- the way each present actualization is always already shaped by prior actualizations that have sedimented as dispositional ground.

The chain runs synchronically from $A_0$ to $A_4$ within any single moment of \textit{pathos}-actualization, but each moment's $A_0$ is itself the having-been-shaped of prior runs. The actualization chain is, in this precise sense, synchronically directional but diachronically saturated. With this temporal structure now specified --- the chain's synchronic directionality (established in \S1.5--\S1.6) and its diachronic temporality (established here) both made explicit --- the analysis can turn, in \S1.8, to the question of emotion's specific causal function within the chain: what work, that is, emotion does at the $M_3 \rightarrow M_4$ transition that \textit{orexis} alone could not.
